% Part: computability
% Chapter: computability-theory
% Section: k-1

\documentclass[../../../include/open-logic-section]{subfiles}

\begin{document}

\olfileid{cmp}{thy}{k1}
\olsection{తగ్గించదగినతకు ఒక ఉదాహరణ}

\olref[ppr]{prop:reduce} యొక్క ఒక అనువర్తనాన్ని పరిశీలిద్దాం.

\begin{prop}
\ollabel{prop:k1}
\[
K_1=\Setabs{e}{\cfind{e}(0)\fdefined}.
\]
అని ఉంచండి. అప్పుడు $K_1$ గణనీయంగా లెక్కించదగినది, కాని గణనీయం
కాదు.
\end{prop}

\begin{proof}
$K_1=\Setabs{e}{\lexists[s][T(e,0,s)]}$ కనుక,
\olref[eqc]{thm:exists-char} ప్రకారం $K_1$ గణనీయంగా లెక్కించదగినది.

$K_1$ గణనీయం కాదని చూపడానికి, $K_0$ దానికి తగ్గించదగినదని చూపుదాం.

\begin{explain}
ఇది కొంచెం చాకచక్యంతో చేయాలి. ఎందుకంటే $K_1$ను వాడి, ఒక నిర్దిష్ట
ఇన్‌పుట్~$0$తో మొదలయ్యే గణనల గురించే ప్రశ్నలు అడగగలం. ఈ రకం
ప్రశ్నలకు సమాధానం చెప్పగల తెలివైన మిత్రుడు మీకు ఉన్నాడనుకోండి
(ఇలాంటి మిత్రులను ``ఒరాకిళ్లు'' అంటారు). అప్పుడు ఎవరో మీ దగ్గరకు వచ్చి
$\tuple{e,x}$ అనేది $K_0$లో ఉందా, అంటే యంత్రం~$e$ ఇన్‌పుట్~$x$పై
ఆగుతుందా అని అడిగారనుకోండి. మీరు చేయగల ఒక పని ఏమిటంటే,
\emph{ఏ} ఇన్‌పుట్ వచ్చినా దాన్ని విస్మరించి, బదులుగా ఇన్‌పుట్~$x$పై
$e$ను నడిపే మరో యంత్రం~$e_x$ను నిర్మించడం. యంత్రం~$e$ ఇన్‌పుట్~$x$పై
ఆగుతుందా అనే ప్రశ్న, యంత్రం~$e_x$ ఇన్‌పుట్~$0$పై (లేదా మరే
ఇన్‌పుట్‌పైనైనా) ఆగుతుందా అనే ప్రశ్నకు స్పష్టంగా తుల్యం. కాబట్టి ఈ
కొత్త యంత్రం~$e_x$ ఇన్‌పుట్~$0$పై ఆగుతుందా అని మీ మిత్రుణ్ణి అడగండి;
మార్చిన ప్రశ్నకు మిత్రుడు ఇచ్చే సమాధానమే అసలు ప్రశ్నకు సమాధానం.
ఇలా $K_0$ నుంచి~$K_1$కు కావలసిన తగ్గింపు లభిస్తుంది.
\end{explain}

సార్వత్రిక పాక్షిక గణనీయ ప్రమేయాన్ని వాడి, 3-స్థాన ప్రమేయం~$f$ను
\[
f(x,y,z)\simeq\cfind{x}(y).
\]
అని నిర్వచించండి. $f$ తన మూడవ ఇన్‌పుట్‌ను పూర్తిగా విస్మరిస్తుందని
గమనించండి. $f=\cfind{e}[3]$ అయ్యే సూచిక~$e$ను ఎంచుకోండి; అప్పుడు
\[
\cfind{e}[3](x,y,z)\simeq\cfind{x}(y).
\]
$s$-$m$-$n$ సిద్ధాంతం ప్రకారం, ప్రతి~$z$కు
\begin{align*}
\cfind{s(e,x,y)}(z) & \simeq \cfind{e}[3](x,y,z) \\
& \simeq \cfind{x}(y).
\end{align*}
అయ్యే ప్రమేయం $s(e,x,y)$ ఒకటి ఉంది.

\begin{explain}
పై అనౌపచారిక వాదన పరంగా, $s(e,x,y)$ అనేది ఏ ఇన్‌పుట్~$z$కైనా దాన్ని
విస్మరించి $\cfind{x}(y)$ను గణించే యంత్రానికి సూచిక.
\end{explain}

ముఖ్యంగా,
\[
\cfind{s(e,x,y)}(0) \fdefined \quad \text{అయితేనే} \quad
\cfind{x}(y) \fdefined.
\]
మరో మాటలో, $\tuple{x,y}\in K_0$ అయితేనే $s(e,x,y)\in K_1$.
కనుక
\[
g(w)=s(e,(w)_0,(w)_1)
\]
అని నిర్వచించిన ప్రమేయం~$g$, $K_0$ను~$K_1$కు తగ్గిస్తుంది.
\end{proof}

\end{document}
